% =========================================================================
% SciPost LaTeX template
% Version 1e (2017-10-31)
%
% Submissions to SciPost Journals should make use of this template.
%
% INSTRUCTIONS: simply look for the `TODO:' tokens and adapt your file.
%
% - please enable line numbers (package: lineno)
% - you should run LaTeX twice in order for the line numbers to appear
% =========================================================================


% TODO: uncomment ONE of the class declarations below
% If you are submitting a paper to SciPost Physics: uncomment next line
\documentclass[preprint]{SciPost}
% If you are submitting a paper to SciPost Physics Lecture Notes: uncomment next line
%\documentclass[submission, LectureNotes]{SciPost}
% If you are submitting a paper to SciPost Physics Proceedings: uncomment next line
%\documentclass[submission, Proceedings]{SciPost}

% % % % %
%SCIPOST RECOMMENDED
\usepackage{doi}
\usepackage{lineno}
%Following definition of 'lightblue' obtained from http://latexcolor.com/
\definecolor{lightblue}{rgb}{0.68, 0.85, 0.9}
\definecolor{airforceblue}{rgb}{0.36, 0.54, 0.66}

%\usepackage[colorlinks=true,citecolor=blue,linkcolor=blue,urlcolor=airforceblue,breaklinks]{hyperref}
%END SCIPOST RECOMMENDED
% % % % %

\usepackage{comment}

\usepackage{amsmath,amssymb,amsfonts,amsthm}
\usepackage{enumerate}
\usepackage{enumitem}
\usepackage{latexsym}
\usepackage{xcolor}

%\usepackage{url}
%\usepackage{breakurl} 

% % % Custom Commands
\newcommand{\bite}{\begin{itemize}}
\newcommand{\eat}{\end{itemize}}
\newcommand{\beq}{\begin{equation}}
\newcommand{\eeq}{\end{equation}}
\newcommand{\rarrow}{\rightarrow}
\newcommand{\bra}{\langle}
\newcommand{\ket}{\rangle}
\newcommand{\beqa}{\begin{align}}
\newcommand{\eeqa}{\end{align}}
\newcommand{\barr}{\begin{array}}
\newcommand{\earr}{\end{array}}
\newcommand{\del}{\partial}
\newcommand{\de}{\mathrm{d}}
%\newcommand{\mu\nu}{{\mu\nu}}
\renewcommand{\th}{\mathrm{th}}
\newcommand{\com}[1]{\begin{itemize}\color{RED}{{#1}}\end{itemize}}
\newcommand{\C}{\mathbb{C}}
\newcommand{\R}{\mathbb{R}}
\newcommand{\btw}[1]{\color{PURPLE}{{#1}}\color{BLACK}}
\newcommand{\cut}[1]{\color{RED}{{#1}}\color{BLACK}}

%text
\newcommand{\ie}{\textit{i.e.}~}
\newcommand{\eg}{\textit{e.g.}~}
\newcommand{\wrt}{\textit{w.r.t.}~}
\newcommand{\etc}{\textit{etc.}~}


\newcommand{\M}{\mathcal{M}}
\newcommand{\N}{\mathcal{N}}
% \newcommand{\H}{\mathcal{H}}

\newcommand{\mb}[1]{\mathbf{#1}}
\newcommand{\mc}[1]{\mathcal{#1}}
\newcommand{\mbb}[1]{\mathbb{#1}}
\newcommand{\mf}[1]{\mathfrak{#1}}

\newcommand{\vect}[1]{\boldsymbol{#1}}
\newcommand{\expect}[1]{\langle #1\rangle}
\newcommand{\innerp}[2]{\langle #1 \vert #2 \rangle}
\newcommand{\fullbra}[1]{\langle #1 \vert}
\newcommand{\fullket}[1]{\vert #1 \rangle}
\newcommand{\supersc}[1]{$^{\textrm{#1}}$}
\newcommand{\subsc}[1]{$_{\textrm{#1}}$}
\newcommand{\sltwoc}{\mathfrak{sl}(2,\mathbb{C})}
% % %

%\keywords{quantum gravity, time reversal symmetry, gauge field, symmetry breaking, tensor networks, arrow of time}


\begin{document}


% TODO: write your article's title here.
% The article title is centered, Large boldface, and should fit in two lines
\begin{center}{\Large \textbf{
Arrow of Time from Spontaneous Symmetry Breaking in Quantum Gravity
}}\end{center}

% TODO: write the author list here. Use initials + surname format.
% Separate subsequent authors by a comma, omit comma at the end of the list.
% Mark the corresponding author with a superscript *.
\begin{center}
Deepak Vaid\textsuperscript{*}
%A. B. Cee\textsuperscript{2},
%G. K. See\textsuperscript{3*}
\end{center}

% TODO: write all affiliations here.
% Format: institute, city, country
\begin{center}
Department of Physics, National Institute of Technology Karnataka (NITK), India \\
%\\
%{\bf 2} Affiliation2
%\\
%{\bf 3} Affiliation2
%\\
% TODO: provide email address of corresponding author
%* CorrespondingAuthor@email.address
* dvaid79@gmail.com
\end{center}

\begin{center}
\today
\end{center}

% For convenience during refereeing: line numbers
%\linenumbers

\section*{Abstract}
{\bf
The question of the origin time's arrow is a major outstanding problem in physics. Here we present a possible mechanism for generation of a cosmological arrow via spontaneous symmetry breaking in a theory of quantum gravity. In \cite{Chen2015Gauging}, Chen and Vishwanath have shown that a local notion of time-reversal symmetry can be encoded into a gauge field living on a tensor network. Using the fact that tensor networks can be used to describe the state space of Loop Quantum Gravity \cite{Qi2013Exact,Han2016Loop} we show that a spontaneous symmetry breaking mechanism operating on such tensor networks can lead to the spontaneous generation of an arrow of time corresponding to the generation of non-zero order parameter for the Chen-Vishwanath gauge field.
}


% TODO: include a table of contents (optional)
% Guideline: if your paper is longer that 6 pages, include a TOC
% To remove the TOC, simply cut the following block
\vspace{10pt}
\noindent\rule{\textwidth}{1pt}
\tableofcontents\thispagestyle{fancy}
\noindent\rule{\textwidth}{1pt}
\vspace{10pt}


The question of ``time'' has perplexed humanity since the dawn of the civilization. Modern physics provides two conflicting roles for time. In quantum mechanics, time is simply an external parameter representing external clocks \wrt which the evolution of a quantum system takes places. In general relativity, on the other hand, time is a dimension on par with the three spatial dimensions and taken together, the four dimensions constitute a dynamical entity known as ``spacetime'' which serves as the arena on which all physics processes take place. However, unified though it may be in a certain geometrical sense with the spatial dimensions, time retains a character distinct from them. The Universe we live in appears to have a distinct \textit{arrow of time}, a sense in which time \emph{flows} from the ``past'' and towards the ``future''. Spatial dimensions do not appear to display such behavior, with all directions being ``equally likely'' - in the sense that the matter distribution in the Universe does not appear to have a preferred spatial orientation on the largest observable scales\footnote{Whether or not the Universe is spatially isotropic on the largest scales is a matter of some debate \cite{Land2005Examination}.}.

\section{Time Reversal Symmetry}



\section{Spin Networks: Microscopic Structure of Spacetime}

\section{Tensor Networks}

\section{Time Reversal Symmetry in Tensor Networks}

\section{Discussion}

\section*{Acknowledgments}

The author would like to thank ...

% TODO: include funding information
\paragraph{Funding information}

The author wishes to acknowledge the support of ...

\bibliography{timesarrow.bib}

\nolinenumbers

\end{document}
